\documentclass[a4paper,11pt]{article}
\usepackage{ctex}
\usepackage{amsmath, amssymb}
\usepackage{booktabs}
\usepackage{float}
\usepackage{geometry}
\usepackage{array}
\geometry{left=2.5cm, right=2.5cm, top=2.5cm, bottom=2.5cm}

\begin{document}

\begin{center}
    \Large\textbf{上机作业 ：用幂法求多项式方程的模最大根}
\end{center}

\section{问题与算法}
题目要求对首一多项式
\[
    f(x)=x^n+a_{n-1}x^{n-1}+\cdots+a_1x+a_0
\]
求模最大的根。教材第 6.2 节给出的幂法迭代格式为
\[
    y_k=Au_{k-1},\qquad
    \mu_k=y_k^{(j)},\quad |y_k^{(j)}|=\|y_k\|_\infty,\qquad
    u_k=\frac{y_k}{\mu_k}.
\]
当矩阵存在唯一的模最大特征值 $\lambda_1$，且初始向量在对应特征子空间上的投影非零时，
$\mu_k$ 收敛到 $\lambda_1$，$u_k$ 收敛到相应特征向量。

为了把多项式求根转化为矩阵特征值问题，程序构造 Frobenius 伴随矩阵
\[
    C=
    \begin{bmatrix}
    0&1&0&\cdots&0\\
    0&0&1&\cdots&0\\
    \vdots&\vdots&\vdots&\ddots&\vdots\\
    0&0&0&\cdots&1\\
    -a_0&-a_1&-a_2&\cdots&-a_{n-1}
    \end{bmatrix}.
\]
直接计算可得
\[
    \det(xI-C)=x^n+a_{n-1}x^{n-1}+\cdots+a_1x+a_0=f(x),
\]
所以 $f(x)=0$ 的根正是 $C$ 的特征值。于是对 $C$ 应用幂法即可得到多项式的模最大根。

程序采用确定性初始向量
\[
    u_0=(1,1+\tfrac1{n-1},\ldots,2)^T
\]
并先按无穷范数归一化。停止准则为
\[
    \max\left\{
    \frac{|\mu_k-\mu_{k-1}|}{\max(|\mu_k|,1)},
    \frac{\|u_k-u_{k-1}\|_\infty}{\max(\|u_k\|_\infty,1)}
    \right\}<10^{-12},
\]
同时检查伴随矩阵特征残差
\[
    \frac{\|Cu_k-\mu_k u_k\|_\infty}
    {\max(\|C\|_\infty\|u_k\|_\infty,1)}.
\]
参考根只用于结果核验，不参与幂法子程序本身。

\section{程序结构}
本次作业的 Python 程序为 \verb|hw7.py|，核心函数如下。
\begin{itemize}
    \item \verb|companion_matrix(coefficients)|：由 $[a_{n-1},\ldots,a_0]$ 构造伴随矩阵。
    \item \verb|power_method_dominant_eigenvalue(A)|：按教材公式实现无穷范数归一化幂法。
    \item \verb|modulus_largest_root(coefficients)|：先构造伴随矩阵，再调用幂法得到模最大根。
    \item \verb|run_all_cases()|：对题目三组多项式运行实验并输出表格、CSV 和 Markdown 结果。
\end{itemize}
程序另外给出两个校验量：一是伴随矩阵的相对特征残差；二是按多项式各项绝对值和缩放后的残差
\[
    \frac{|f(\hat x)|}
    {|\hat x|^n+|a_{n-1}||\hat x|^{n-1}+\cdots+|a_1||\hat x|+|a_0|}.
\]
后者用于避免高次多项式在大根处直接代入时，由巨大中间量相互抵消造成的绝对残差误判。

\section{数值结果}
表 \ref{tab:results} 给出三道题的幂法计算结果。表中 ``scaled residual'' 是上节定义的缩放多项式残差，
``matrix rel residual'' 是伴随矩阵相对特征残差。

\begin{table}[H]
\centering
\scriptsize
\caption{幂法求多项式模最大根的结果}
\label{tab:results}
\setlength{\tabcolsep}{3pt}
\begin{tabular}{cccrrrr}
\toprule
题号 & 次数 & 幂法结果 & 迭代次数 & scaled residual & matrix rel residual & 参考模最大根 \\
\midrule
(i) & 3 & $-3.00000000000069$ & 65  & $2.053\times10^{-13}$ & $1.030\times10^{-13}$ & $-3$ \\
(ii) & 3 & $1.87938524157102$ & 136 & $4.548\times10^{-13}$ & $3.607\times10^{-13}$ & $1.87938524157182$ \\
(iii) & 8 & $-99.9999999999989$ & 15 & $5.294\times10^{-15}$ & $5.352\times10^{-18}$ & $-100$ \\
\bottomrule
\end{tabular}
\end{table}

\subsection{第 (i) 题}
\[
    x^3+x^2-5x+3=(x-1)^2(x+3).
\]
全部根为 $1,1,-3$，因此模最大根唯一为 $-3$。幂法 65 次后得到
$\hat x=-3.00000000000069$，与理论根一致。

\subsection{第 (ii) 题}
\[
    x^3-3x-1=0.
\]
令 $x=2\cos\theta$，则 $x^3-3x=2\cos3\theta$，所以
\[
    \cos3\theta=\frac12.
\]
三根可写为
\[
    2\cos\frac{\pi}{9},\quad
    2\cos\frac{5\pi}{9},\quad
    2\cos\frac{7\pi}{9},
\]
数值上为
\[
    1.87938524157182,\quad -0.347296355333861,\quad -1.53208888623796.
\]
模最大根为 $2\cos(\pi/9)\approx1.87938524157182$。幂法得到
$1.87938524157102$，相对特征残差约为 $3.607\times10^{-13}$。

\subsection{第 (iii) 题}
题目多项式为
\[
\begin{aligned}
f(x)=&x^8+101x^7+208.01x^6+10891.01x^5+9802.08x^4\\
&+79108.9x^3-99902x^2+790x-1000 .
\end{aligned}
\]
为了核验结果，可将其因式分解为
\[
    f(x)=\frac{1}{100}(x-1)(x+100)(x^2+100)(100x^2+1)(x^2+2x+10).
\]
因此全部根为
\[
    -100,\quad \pm10i,\quad -1\pm3i,\quad 1,\quad \pm0.1i,
\]
模最大根唯一为 $-100$。幂法只需 15 次迭代即得到
$\hat x=-99.9999999999989$。虽然把这个近似值直接代入八次多项式会产生较大的未缩放绝对残差，
但这是因为 $x\approx -100$ 时各项量级可达 $10^{16}$ 并发生严重消去；缩放残差为
$5.294\times10^{-15}$，伴随矩阵相对特征残差为 $5.352\times10^{-18}$，说明所得根可靠。

\section{结论}
本次实验按教材第 6.2 节的幂法格式实现了一个通用子程序：
先把首一多项式转化为 Frobenius 伴随矩阵，再用无穷范数归一化幂法求该矩阵的模最大特征值。
三道题的模最大根分别为
\[
    -3,\qquad 1.87938524157182,\qquad -100.
\]
实验结果与因式分解或三角表示得到的参考根一致。第 (iii) 题也说明，在高次多项式且根的模很大时，
应同时观察缩放残差和伴随矩阵特征残差，不能只用未缩放的 $|f(\hat x)|$ 判断计算质量。

\end{document}
